\chapter{Anforderungsanalyse}
\textit{„Eine Anforderungsspezifikation bzw. ein Anforderungsdokument ist eine systematisch dargestellte Sammlung von Anforderungen (typischerweise für ein System oder eine Komponente), die vorgegebenen Kriterien genügt.“} [PoRu15, S. 35, Definition 4-1] %% ToDo bib.tex

% TODO: bib.tex
Die Anforderungen können nach [PoRu15] von verschiedene Perspektiven betrachtet werden. Zur Beschreibung der Anforderungen des zu entwickelnden Algorithmus eignen sich die Strukturperspektive und die Funktionsperspektive. In der Betrachtung der Funktionsperspektive kann zwischen den funktionalen Anforderungen, Qualitätsanforderungen und Randbedingung unterschieden werden.

Die Anforderungen basieren auf der Absprache mit Petrit Vuthi vom CC4E, der den Algorithmus im Rahmen seiner Doktorarbeit einsetzen wird.

\section{Strukturperspektive}
(Ein- und Ausgabedaten)
\begin{enumerate}[label=Req-S-\arabic*:]
    \item Der zu betrachtende Kartenausschnitt (Region) muss im GeoJSON Format eingelesen werden können.
    \item Das Koordinatenreferenzsystem muss EPSG:3857 sein.
    \item Die Gebäude müssen dabei als zweidimensionaler Gebäudegrundrisse (LoD0) dargestellt werden und physische Gebäude darstellen (keine Haushalte).
    \item Es müssen nur Regionen mit einer maximalen Anzahl von 100000 Häusern betrachtet werden können.
    \item Die maximale räumliche Ausdehnung der Region darf 10 km in Ost-West-Richtung und 10 km in Nord-Süd-Richtung nicht überschreiten.
    \item Es müssen Last- und Erzeugungsprofile eingelesen werden können. Diese repräsentieren den generierten und verbrauchten Strom pro abstrakter Zeiteinheit des Hauses. In dem GeoJSON sollen diese unter den Properties Objekt der Häuser in „Generation“ und „Last“ Objekte geschrieben sein.
    \item Die Last- und Erzeugungsprofile sollen aus einer Liste von kWh Werten bestehen, die als Gleitkommazahl kodiert sind. Jedes Haus soll beide Objekte beinhalten und die Listen müssen bei jedem Haus gleich viele Elemente beinhalten.
    \item Diese Werte müssen positiv und vollständig sein.
    \item Die Last- und Erzeugungsprofile können maximal 8784 (Stunden in einem Schaltjahr) Zeitschritte pro Haus lang sein. 
    \item Die Last- und Erzeugungsprofile können maximal Werte von 1000 kWh pro Zeitschritt pro Haus haben. 
    \item Die Ausgabe des Algorithmus soll eine Liste von Label werten sein. Diese sind in derselben Reihenfolge anzuordnen wie die Häuser in dem Eingangs-GeoJSON. 
\end{enumerate}

\section{Funktionale Anforderungen}
\begin{enumerate}[label=Req-F-\arabic*:]
    \item Der Algorithmus soll die Häuser aus dem GeoJSON zu Gruppen zusammenführen.
    \item Diese Gruppen sollen möglichst gut aufgeteilt werden.
    \item Wie gut eine Gruppe aufgeteilt wurde bestimmt eine Bewertungsfunktion (Region Score). Diese wird durch die Formel im Kapitel „Technische Analyse“ beschrieben.
    \item Die Anzahl der Gruppen muss im Vorhinein nicht festgelegt sein.
    \item Die maximale Gruppengröße muss festgelegt werden können.
    \item Ein Haus kann nur einer Gruppe zugeordnet werden.
    \item Der Algorithmus soll unter Eingabe des gleichen Seeds deterministisch sein.
\end{enumerate}

\section{Qualitätsanforderungen}
\begin{enumerate}[label=Req-Q-\arabic*:]
    \item Die Laufzeit soll in \(O(n)\) liegen, wobei \(n\) der Anzahl der betrachteten Häuser entspricht.
    \item Die Laufzeit darf auf einem System mit einem Intel Core i7-XXXX, XX GB RAM und Windows/Linux bei der Verarbeitung von XX Häusern maximal XX Minuten betragen.
\end{enumerate}

\section{Randbedingungen}
\begin{enumerate}[label=Req-R-\arabic*:]
    \item Wenn die Strukturanforderungen der Eingabe nicht eingehalten werden können, bricht der Algorithmus mit einem Fehler die Ausführung ab.
\end{enumerate}
